% Capítulo gerado a partir de zeros_de_funcoes.tex (arquivos usados neste capítulo)
\def\arqBisA{metodoBisseccao_2026_corrigido.ggb}
\def\arqBisB{metodoBisseccao_2026_correto_and.ggb}
\def\arqPosA{metodoposicaofalsa_2026_correto.ggb}
\def\arqPosB{metodoposicaofalsa_2026_correto_and.ggb}
\def\arqFixo{metodopontofixoJAVASCRIPT_2026_final_corrigido.ggb}
\def\arqFixoC{metodopontofixoJAVASCRIPT_2026_CORRIGIDO_ex02_CONVERGINDO.ggb}
\def\arqFixoD{metodopontofixoJAVASCRIPT_2026_CORRIGIDO_ex02_DIVERGINDO.ggb}
\def\arqNewton{metodoNewtonRaphson_2026_corrigido.ggb}
\def\arqSec{metododasecanteJAVASCRIPT_2026_CORRIGIDO.ggb}
\def\pdfLacos{material_lacos_fluxograma.pdf}
\chapter{Zeros de funções}\label{cap:zeros}

\begin{tcolorbox}[colback=gray!4,colframe=gray!50,title={Como usar os aplicativos e os textos deste capítulo},fonttitle=\bfseries]
Os quadros laranja \emph{Atividade no GeoGebra} trazem o link \textbf{[abrir o aplicativo]}.
São nove aplicativos:
\begin{itemize}[nosep]
  \item bisseção: \ggbabrir{\arqBisA}{versão 1, em $[1{,}6;\,1{,}8]$} e \ggbabrir{\arqBisB}{versão 2, em $[1{,}2;\,2]$};
  \item posição falsa: \ggbabrir{\arqPosA}{versão 1} e \ggbabrir{\arqPosB}{versão 2, com tolerâncias};
  \item ponto fixo: \ggbabrir{\arqFixo}{$\varphi(x)=\frac{2(1+x)}{x+2}$},
        \ggbabrir{\arqFixoC}{exemplo que converge} e \ggbabrir{\arqFixoD}{exemplo que diverge};
  \item \ggbabrir{\arqNewton}{Newton-Raphson} e \ggbabrir{\arqSec}{secante}.
\end{itemize}
Em todos eles há três botões: \textbf{INICIAR} (restaura os dados iniciais), \textbf{Xk}
(faz uma iteração por clique, mostrando no gráfico a construção do método) e
\textbf{JAVASCRIPT} (executa o laço completo e mostra a tabela das iterações numa janela).

Os quadros verdes \emph{Leitura complementar} abrem dois textos de apoio:
\pdfabrir{\pdfLacos}{laços de repetição e fluxogramas}.

Para que os links funcionem, mantenha este PDF e as pastas \texttt{ggb} e \texttt{complementos}
juntos e abra o PDF num leitor que permita abrir arquivos externos (Adobe Acrobat Reader, Okular,
SumatraPDF, Foxit). Os leitores de PDF dos navegadores costumam bloquear esses links; nesse caso,
abra os arquivos diretamente das pastas.
\end{tcolorbox}

% ==================================================================
\section{O problema}

Queremos encontrar $\bar x$ tal que $f(\bar x)=0$, isto é, um \emph{zero} (ou \emph{raiz}) de
$f$. Para polinômios de grau até 4 há fórmulas, mas para equações como
\[
x\ln x-1=0,\qquad x^5-x^3+x^2-3=0,\qquad \cos x=x
\]
não há expressão fechada para a solução. Os métodos numéricos produzem uma sequência de
aproximações $x_0,x_1,x_2,\dots$ que converge para $\bar x$, e param quando a aproximação é
boa o bastante.

O trabalho tem duas fases:
\begin{enumerate}[nosep]
  \item \textbf{isolamento}: encontrar um intervalo $[a,b]$ que contenha uma única raiz
        (ou uma boa aproximação inicial);
  \item \textbf{refinamento}: melhorar a aproximação até a precisão desejada.
\end{enumerate}

Ao longo do capítulo, usaremos a função $f(x)=x^2-2$, cujo zero positivo é
$\sqrt2\approx1{,}414214$, para comparar os métodos. Como sabemos a resposta, fica fácil
observar a velocidade de cada um.

% ==================================================================
\section{Isolamento das raízes}

\begin{teorema}[Bolzano]
Se $f$ é contínua em $[a,b]$ e $f(a)\,f(b)<0$, então existe pelo menos um
$\bar x\in(a,b)$ com $f(\bar x)=0$. Se, além disso, $f'$ existe e não muda de sinal em
$(a,b)$, a raiz é única.
\end{teorema}

\begin{exemplo}[{: tabela de sinais}]
Para $f(x)=x^3-9x+3$:
\begin{center}
\begin{tabular}{c|ccccccccc}
$x$ & $-4$ & $-3$ & $-2$ & $-1$ & $0$ & $1$ & $2$ & $3$ & $4$\\ \hline
$f(x)$ & $-25$ & $3$ & $13$ & $11$ & $3$ & $-5$ & $-7$ & $3$ & $31$\\
sinal & $-$ & $+$ & $+$ & $+$ & $+$ & $-$ & $-$ & $+$ & $+$
\end{tabular}
\end{center}
Há mudanças de sinal em $[-4,-3]$, $[0,1]$ e $[2,3]$. Como $f$ tem grau 3, essas são as três raízes.
\end{exemplo}

\begin{exemplo}[{: a função dos aplicativos de bisseção}]
Para $f(x)=x\ln x-1$ (com $x>0$), temos $f(1)=-1<0$ e $f(2)=2\ln2-1\approx0{,}386>0$.
Como $f'(x)=\ln x+1>0$ em $[1,2]$, há uma única raiz nesse intervalo:
$\bar x\approx1{,}763223$. Graficamente, ela é a abscissa do ponto em que $\ln x=\frac1x$.
\end{exemplo}

\begin{observacao}
Mudança de sinal é condição suficiente, não necessária: uma raiz de multiplicidade par não
troca o sinal de $f$. Por exemplo, $f(x)=x^3-5x^2+8x-4=(x-1)(x-2)^2$ (a função dos aplicativos
de posição falsa) tem uma raiz dupla em $x=2$, onde o gráfico só tangencia o eixo. Os métodos de
intervalo (bisseção e posição falsa) não conseguem ``cercar'' essa raiz.
\end{observacao}

% ==================================================================
\section{Critérios de parada e laços de repetição}

Todos os métodos deste capítulo são \emph{iterativos}: repetem o mesmo cálculo até que um
critério de parada seja satisfeito. Dada uma tolerância $\varepsilon$, os critérios mais usados são
\[
\text{(i) } |x_k-x_{k-1}|<\varepsilon,\qquad
\text{(ii) } |f(x_k)|<\varepsilon,\qquad
\text{(iii) } \frac{|x_k-x_{k-1}|}{|x_k|}<\varepsilon,
\]
\[
\text{(iv) } |b_k-a_k|<\varepsilon\ \text{(métodos de intervalo)},
\]
sempre acompanhados de um número máximo de iterações, para evitar um laço infinito quando o
método não converge.

Nenhum critério isolado é perfeito: $|f(x_k)|$ pode ser pequeno longe da raiz (função muito
``achatada'') e $|x_k-x_{k-1}|$ pode ser pequeno sem que se esteja perto dela (convergência muito
lenta). Por isso, é comum combinar dois critérios, e a escolha entre \emph{e} e \emph{ou} importa:
\begin{itemize}[nosep]
  \item \textbf{continuar enquanto} $|b-a|>\varepsilon_1$ \textbf{ou} $|f(x_k)|>\varepsilon_2$:
        só para quando \emph{os dois} critérios forem satisfeitos (mais exigente). É o que fazem os
        aplicativos de bisseção, com as tolerâncias $tol_1$ e $tol_2$;
  \item \textbf{continuar enquanto} $|x_k-x_{k-1}|>\varepsilon$ \textbf{e} $|f(x_k)|>\varepsilon$:
        para quando \emph{um} deles for satisfeito (menos exigente). É o que fazem os botões
        JAVASCRIPT dos aplicativos de posição falsa, secante, Newton e ponto fixo.
\end{itemize}

\begin{figure}[h]
\centering
\begin{tikzpicture}[node distance=7mm and 10mm,
  bloco/.style={draw,rounded corners,fill=blue!6,minimum width=46mm,minimum height=7mm,align=center,font=\small},
  decisao/.style={draw,diamond,aspect=2.4,fill=orange!10,align=center,font=\small,inner sep=1pt},
  seta/.style={-{Stealth[length=2.5mm]}}]
\node[bloco] (ini) {Início: dados $f$, $x_0$ (ou $[a,b]$), $\varepsilon$, $N$};
\node[bloco,below=of ini] (k0) {$k\leftarrow0$};
\node[decisao,below=of k0] (teste) {critério de parada\\satisfeito ou $k\ge N$?};
\node[bloco,below=9mm of teste] (passo) {calcula $x_{k+1}$ pela fórmula do método};
\node[bloco,below=of passo] (inc) {$k\leftarrow k+1$};
\node[bloco,right=16mm of teste] (fim) {Fim: $\bar x\approx x_k$};
\draw[seta] (ini)--(k0); \draw[seta] (k0)--(teste);
\draw[seta] (teste)--node[right,font=\small]{não}(passo);
\draw[seta] (passo)--(inc);
\draw[seta] (inc.west)--++(-12mm,0)|-(teste.west);
\draw[seta] (teste)--node[above,font=\small]{sim}(fim);
\end{tikzpicture}
\caption{Fluxograma comum a todos os métodos iterativos: um laço \texttt{while}.}
\label{fig:fluxo}
\end{figure}

Em linguagem de programação, esse fluxograma é um laço \texttt{while}. O botão JAVASCRIPT dos
aplicativos de bisseção, por exemplo, executa (resumidamente):
\begin{lstlisting}[language=Java]
var k = 0;  var xk = (a + b) / 2;
while (F(xk) !== 0 && (Math.abs(b - a) > tol1 || Math.abs(F(xk)) > tol2) && k < 200) {
    if (F(a) * F(xk) > 0) { a = xk; } else { b = xk; }   // escolhe o novo intervalo
    k++;
    xk = (a + b) / 2;                                    // novo ponto medio
}
\end{lstlisting}
A condição \texttt{k < 200} é o número máximo de iterações, que impede um laço infinito.

\begin{leitura}{laços de repetição e fluxogramas (4 páginas)}{\pdfLacos}
Introdução aos laços \texttt{for} e \texttt{while}, com exemplos em Python e fluxogramas:
variável de controle, condição lógica, contador e acumulador, laço infinito, \texttt{break} e
\texttt{continue}. É a base para ler e escrever o código dos botões JAVASCRIPT dos aplicativos.
\end{leitura}

% ==================================================================
\section{Método da bisseção}

Partindo de $[a_0,b_0]$ com $f(a_0)f(b_0)<0$, calculamos o ponto médio
\[
x_k=\frac{a_k+b_k}{2}
\]
e ficamos com a metade em que há troca de sinal: se $f(a_k)f(x_k)<0$, então
$[a_{k+1},b_{k+1}]=[a_k,x_k]$; caso contrário, $[a_{k+1},b_{k+1}]=[x_k,b_k]$.

\begin{proposicao}[Estimativa do erro]
Como o intervalo é dividido ao meio a cada passo,
$|\bar x-x_k|\le\dfrac{b_k-a_k}{2}=\dfrac{b_0-a_0}{2^{k+1}}$. Para garantir erro menor que
$\varepsilon$ com o critério $b_k-a_k<\varepsilon$, bastam
\[
k>\log_2\frac{b_0-a_0}{\varepsilon}
\]
iterações, número que se conhece \emph{antes} de começar.
\end{proposicao}

\begin{exemplo}[{: três iterações para $f(x)=x^2-2$ em $[1,2]$}]
$f(1)=-1<0$ e $f(2)=2>0$.
\begin{center}
\begin{tabular}{cccccc}
\toprule
$k$ & $a_k$ & $b_k$ & $x_k=\frac{a_k+b_k}{2}$ & $f(x_k)$ & novo intervalo\\
\midrule
1 & 1 & 2 & 1{,}5 & $\phantom{-}0{,}25$ & $[1;\,1{,}5]$\\
2 & 1 & 1{,}5 & 1{,}25 & $-0{,}4375$ & $[1{,}25;\,1{,}5]$\\
3 & 1{,}25 & 1{,}5 & 1{,}375 & $-0{,}109375$ & $[1{,}375;\,1{,}5]$\\
\bottomrule
\end{tabular}
\end{center}
Em cada linha, o novo intervalo é escolhido pelo sinal de $f(x_k)$. Para obter
$b_k-a_k<10^{-3}$, seriam necessárias $\log_2(1/10^{-3})\approx9{,}97$, isto é, 10 iterações.
\end{exemplo}

\begin{figure}[h]
\centering
\begin{tikzpicture}
\begin{axis}[width=11cm,height=6cm,axis lines=middle,xmin=0.8,xmax=2.2,ymin=-1.3,ymax=2.3,
  xtick={1,1.25,1.375,1.5,2},xticklabels={$1$,$1{,}25$,$1{,}375$,$1{,}5$,$2$},
  ytick={-1,1,2},grid=both,grid style={gray!20},xlabel={$x$},ylabel={$y$}]
\addplot[domain=0.85:2.1,samples=100,very thick,azul]{x^2-2};
\addplot[only marks,mark=*,vermelho] coordinates {(1.5,0)(1.25,0)(1.375,0)};
\addplot[only marks,mark=*,azul] coordinates {(1,-1)(2,2)(1.5,0.25)(1.25,-0.4375)(1.375,-0.109375)};
\draw[thick,laranja,|-|] (axis cs:1,-1.1)--(axis cs:2,-1.1) node[right,font=\scriptsize]{$[a_0,b_0]$};
\draw[thick,laranja,|-|] (axis cs:1,-0.85)--(axis cs:1.5,-0.85);
\draw[thick,laranja,|-|] (axis cs:1.25,-0.6)--(axis cs:1.5,-0.6);
\draw[thick,laranja,|-|] (axis cs:1.375,-0.35)--(axis cs:1.5,-0.35);
\end{axis}
\end{tikzpicture}
\caption{Bisseção para $f(x)=x^2-2$: a cada passo o intervalo (em laranja) cai pela metade.}
\end{figure}

\begin{atividade}{bisseção (versão 1)}{\arqBisA}
A função é $f(x)=x\ln x-1$. O botão INICIAR usa $[a,b]=[1{,}6;\,1{,}8]$ e
$tol_1=tol_2=0{,}001$, que podem ser alteradas nas caixas $tol_1$ e $tol_2$. O botão Xk só avança
enquanto $f(a)f(b)<0$ e ($|b-a|>tol_1$ ou $|f(x_k)|>tol_2$).
\begin{enumerate}[nosep]
  \item Clique em INICIAR e depois em Xk repetidas vezes, anotando $a_k$, $b_k$, $x_k$ e
        $f(x_k)$. Confira: $x_0=1{,}7$; $x_1=1{,}75$; $x_2=1{,}775$; $x_3=1{,}7625$.
  \item Em que iteração o botão deixa de avançar? (Esperado: $k=8$, com $[a,b]=[1{,}7625;\,1{,}763281]$,
        $|b-a|\approx0{,}00078$ e $|f(x_k)|\approx0{,}00052$.) Compare com a estimativa
        $\log_2(0{,}2/0{,}001)\approx7{,}6$.
  \item Use o botão JAVASCRIPT para ver todas as iterações de uma vez. Ele lê $a$, $b$, $tol_1$ e
        $tol_2$ do aplicativo; teste com $tol_1=tol_2=10^{-6}$.
  \item Escolha com os controles deslizantes um intervalo sem troca de sinal, por exemplo
        $[2;3]$. O que acontece?
\end{enumerate}
\end{atividade}

\begin{atividade}{bisseção (versão 2, intervalo maior)}{\arqBisB}
Mesma função, mas o botão INICIAR usa $[a,b]=[1{,}2;\,2]$, e o texto mostra também os dados da
iteração inicial ($a_0$, $b_0$, $x_0$). Os dois arquivos de bisseção trazem ainda uma tabela na
Planilha (menu Exibir $\to$ Planilha) que executa o algoritmo célula a célula, com as colunas
$a$, $b$, $x_k$, $f(a)$, $f(x_k)$, $f(a)f(x_k)$, $a_k$, $b_k$ e $|b_k-a_k|$.
\begin{enumerate}[nosep]
  \item Quantas iterações são necessárias agora? (Esperado: $k=10$.) Por que são duas a mais que na
        versão 1? (Sugestão: compare $b_0-a_0$ nos dois casos.)
  \item Na Planilha, a primeira linha da tabela usa as células $B22$ ($a$) e $C22$ ($b$). Digite
        nelas $1$ e $2$ e acompanhe linha a linha: como cada linha obtém $a_k$ e $b_k$ a partir da
        anterior? A coluna $|b_k-a_k|$ é sempre dividida por 2?
\end{enumerate}
\end{atividade}

% ==================================================================
\section{Método da posição falsa}

Em vez do ponto médio, usamos o ponto em que a reta que liga $(a,f(a))$ a $(b,f(b))$ corta o
eixo $x$:
\[
x_k=\frac{a_k\,f(b_k)-b_k\,f(a_k)}{f(b_k)-f(a_k)}.
\]
A escolha do novo intervalo é igual à da bisseção. Quando $|f(a)|$ e $|f(b)|$ são muito
diferentes, $x_k$ tende a ficar mais perto da raiz do que o ponto médio.

\begin{exemplo}[{: três iterações para $f(x)=x^2-2$ em $[1,2]$}]
\begin{center}
\begin{tabular}{cccccc}
\toprule
$k$ & $a_k$ & $b_k$ & $x_k$ & $f(x_k)$ & novo intervalo\\
\midrule
1 & 1 & 2 & $4/3\approx1{,}333333$ & $-0{,}222222$ & $[1{,}333333;\,2]$\\
2 & 1{,}333333 & 2 & 1{,}400000 & $-0{,}040000$ & $[1{,}4;\,2]$\\
3 & 1{,}4 & 2 & 1{,}411765 & $-0{,}006920$ & $[1{,}411765;\,2]$\\
\bottomrule
\end{tabular}
\end{center}
Por exemplo, $x_1=\frac{1\cdot2-2\cdot(-1)}{2-(-1)}=\frac43$. Repare que a extremidade $b=2$
\emph{nunca muda}: como $f$ é convexa, todas as aproximações ficam do mesmo lado da raiz. Assim,
$b_k-a_k$ não tende a zero, e o critério de parada deve usar $|f(x_k)|$ ou $|x_k-x_{k-1}|$.
\end{exemplo}

\begin{atividade}{posição falsa (versão 1)}{\arqPosA}
A função é $f(x)=x^3-5x^2+8x-4=(x-1)(x-2)^2$, com $[a,b]=[0{,}5;\,1{,}7]$ no botão INICIAR. O
gráfico mostra a reta secante $g$ que define cada $x_k$.
\begin{enumerate}[nosep]
  \item Confira $x_1\approx1{,}636364$, $x_2\approx1{,}557282$, $x_3\approx1{,}463714$.
  \item Qual extremidade fica fixa? Avance 10 iterações: $x_{10}\approx1{,}022$, ainda longe de
        $\bar x=1$. Por que a convergência é tão lenta aqui? (Observe o formato do gráfico perto de
        $x=1$ e o valor de $f(a)$.)
  \item Tente um intervalo que contenha $x=2$, como $[1{,}5;\,2{,}5]$. O aplicativo avisa
        ``$f(a)f(b)>0$''. Explique usando a raiz dupla.
  \item O botão JAVASCRIPT resolve outro exemplo, $f(x)=x^2-2$ em $[0{,}5;\,2]$ com $\varepsilon=10^{-4}$.
        Compare a velocidade com a da função do gráfico.
\end{enumerate}
\end{atividade}

\begin{atividade}{posição falsa (versão 2, com tolerâncias)}{\arqPosB}
Mesma função, agora com $[a,b]=[0{,}2;\,1{,}5]$ e as tolerâncias $tol_1=tol_2=0{,}001$; o texto
mostra $a_0$, $b_0$, $x_0$ e os erros a cada passo.
\begin{enumerate}[nosep]
  \item Faça 5 iterações e confira $x_5\approx1{,}199928$ (com $a=0{,}2$ fixo).
  \item Compare $|b_k-a_k|$ e $|f(x_k)|$ ao longo das iterações. Qual dos dois critérios seria
        atingido primeiro?
\end{enumerate}
\end{atividade}

% ==================================================================
\section{Método do ponto fixo}

Reescrevemos $f(x)=0$ numa forma equivalente $x=\varphi(x)$ e iteramos
\[
x_{k+1}=\varphi(x_k).
\]
Um ponto com $\bar x=\varphi(\bar x)$ é um \emph{ponto fixo} de $\varphi$ e, portanto, um zero
de $f$. Graficamente, é a interseção de $y=\varphi(x)$ com a reta $y=x$, e as iterações formam o
diagrama em ``teia'' ou ``escada'' (Figura~\ref{fig:teia}).

\begin{teorema}[Convergência do ponto fixo]
Se $\varphi$ e $\varphi'$ são contínuas num intervalo $I$ que contém $\bar x$, $\varphi(I)\subset I$
e $|\varphi'(x)|\le L<1$ para todo $x\in I$, então, para qualquer $x_0\in I$, a sequência
$x_{k+1}=\varphi(x_k)$ converge para o único ponto fixo $\bar x$ em $I$, e
$|x_{k+1}-\bar x|\le L\,|x_k-\bar x|$.
Se $|\varphi'(\bar x)|>1$, as iterações se afastam de $\bar x$.
\end{teorema}

Se $\varphi'>0$ perto da raiz, as aproximações se aproximam de um só lado (escada); se
$\varphi'<0$, elas alternam em torno da raiz (teia). Quanto menor $|\varphi'(\bar x)|$, mais rápida
a convergência.

\begin{exemplo}[{: $f(x)=x^2-2$ com $\varphi(x)=\frac{x+2}{x+1}$ e $x_0=1$}]
De $x^2=2$ vem $x^2+x=x+2$, isto é, $x=\frac{x+2}{x+1}$. Temos
$\varphi'(x)=-\frac{1}{(x+1)^2}$, e em $[1,2]$ vale $|\varphi'(x)|\le\frac14<1$; além disso,
$\varphi([1,2])=[\frac43,\frac32]\subset[1,2]$. Logo, há convergência.
\begin{center}
\begin{tabular}{ccccc}
\toprule
$k$ & $x_k$ & $x_{k+1}=\varphi(x_k)$ & $f(x_{k+1})$ & $|x_{k+1}-x_k|$\\
\midrule
0 & 1 & 1{,}5 & $\phantom{-}0{,}25$ & 0{,}5\\
1 & 1{,}5 & 1{,}4 & $-0{,}04$ & 0{,}1\\
2 & 1{,}4 & $17/12\approx1{,}416667$ & $\phantom{-}0{,}006944$ & 0{,}016667\\
\bottomrule
\end{tabular}
\end{center}
Como $\varphi'<0$, os iterados alternam em torno de $\sqrt2$, e o erro cai cerca de
$|\varphi'(\sqrt2)|\approx0{,}17$ a cada passo.

\textbf{Nem toda $\varphi$ serve.} A mesma equação pode ser escrita como $x=\frac2x$. Agora
$|\varphi'(\sqrt2)|=1$, e com $x_0=1$ obtemos $2,1,2,1,\dots$: a sequência nunca se aproxima de
$\sqrt2$.
\end{exemplo}

\begin{figure}[h]
\centering
\begin{tikzpicture}
\begin{axis}[width=8cm,height=7cm,axis lines=middle,xmin=0.9,xmax=1.6,ymin=0.9,ymax=1.6,
  grid=both,grid style={gray!20},xlabel={$x$},ylabel={$y$},
  legend style={at={(0.02,0.98)},anchor=north west,font=\small}]
\addplot[domain=0.9:1.6,thick,gray]{x};
\addplot[domain=0.9:1.6,very thick,azul,samples=60]{(x+2)/(x+1)};
\legend{$y=x$,$y=\varphi(x)$}
\addplot[thick,vermelho,-{Stealth[length=2mm]}] coordinates {(1,0.9)(1,1.5)(1.5,1.5)(1.5,1.4)(1.4,1.4)(1.4,1.416667)(1.416667,1.416667)};
\addplot[only marks,mark=*,mark size=1.5pt] coordinates {(1,0.9)(1.5,0.9)(1.4,0.9)};
\end{axis}
\end{tikzpicture}
\caption{Diagrama em teia para $\varphi(x)=\frac{x+2}{x+1}$ a partir de $x_0=1$:
os iterados alternam em torno de $\sqrt2$.}
\label{fig:teia}
\end{figure}

\begin{atividade}{ponto fixo com $\varphi(x)=\frac{2(1+x)}{x+2}$}{\arqFixo}
Aqui $f(x)=x^2-2$ e $\varphi(x)=\frac{2(1+x)}{x+2}$ (verifique que $x=\varphi(x)\iff x^2=2$). O
botão INICIAR usa $[a,b]=[1;3]$ e $x_0=b-0{,}3=2{,}7$, e o gráfico mostra a escada (vetores),
$\varphi$ e $|\varphi'|$. O texto indica ``Continue'' ou ``Pare'' conforme os critérios
$|x_{k+1}-x_k|>10^{-5}$ e $|f(x_{k+1})|>10^{-4}$.
\begin{enumerate}[nosep]
  \item Confira $x_1\approx1{,}574468$, $x_2\approx1{,}440476$, $x_3\approx1{,}418685$.
  \item Calcule $\varphi'(x)=\frac{2}{(x+2)^2}$. Por que a convergência é monótona (escada) e não
        alternada? Quanto vale $|\varphi'(\sqrt2)|$?
  \item Use os botões ``Apaga/Mostrar o Graf $f(x)$ e $\varphi'(x)$'' e verifique que
        $|\varphi'(x)|<1$ em todo o intervalo $[1,3]$.
  \item O botão JAVASCRIPT resolve outro exemplo: $x\ln x-1=0$ com $\varphi(x)=e^{1/x}$ e $x_0=1{,}3$.
        Mostre que $\varphi'(\bar x)=-1/\bar x\approx-0{,}57$ e explique por que essa convergência é
        alternada e bem mais lenta.
\end{enumerate}
\end{atividade}

\begin{atividade}{ponto fixo que converge e ponto fixo que diverge}{\arqFixoC}
Os dois aplicativos tratam $f(x)=x^2-x-2=(x-2)(x+1)$, com $x_0=2{,}7$, mas com funções de
iteração diferentes:
\begin{itemize}[nosep]
  \item \ggbabrir{\arqFixoC}{versão que converge}: $\varphi(x)=x-\dfrac{x^2-x-2}{2x-1}$;
  \item \ggbabrir{\arqFixoD}{versão que diverge}: $\varphi(x)=x^2-2$.
\end{itemize}
\begin{enumerate}[nosep]
  \item Na versão que converge, confira $x_1\approx2{,}111364$, $x_2\approx2{,}003848$,
        $x_3\approx2{,}000005$. Que método é esse? (Compare $\varphi$ com a fórmula de Newton.)
  \item Na versão que diverge, confira $x_1=5{,}29$ e $x_2\approx25{,}98$. Calcule $\varphi'(2)=4$
        e explique a divergência pelo teorema.
  \item Nas duas versões, ligue o gráfico de $|\varphi'|$ e compare-o com a reta $y=1$ perto de $x=2$.
\end{enumerate}
\small\textbf{Arquivo da versão que diverge:}\\ \ggbabrir{\arqFixoD}{\texttt{ggb/\arqtxt{\arqFixoD}}}
\end{atividade}

% ==================================================================
\section{Método de Newton-Raphson}

Aproximamos $f$ pela reta tangente em $x_k$ e tomamos o ponto em que ela corta o eixo $x$:
\[
x_{k+1}=x_k-\frac{f(x_k)}{f'(x_k)},\qquad f'(x_k)\neq0.
\]
É o método do ponto fixo com $\varphi(x)=x-\frac{f(x)}{f'(x)}$. Como $\varphi'(\bar x)=0$ quando
$f'(\bar x)\neq0$, a convergência perto de uma raiz simples é \emph{quadrática}: grosso modo, o
número de algarismos corretos dobra a cada iteração.

\begin{exemplo}[{: três iterações para $f(x)=x^2-2$ com $x_0=2$}]
$f'(x)=2x$, e então $x_{k+1}=x_k-\frac{x_k^2-2}{2x_k}$.
\begin{center}
\begin{tabular}{cccccc}
\toprule
$k$ & $x_k$ & $f(x_k)$ & $f'(x_k)$ & $x_{k+1}$ & $f(x_{k+1})$\\
\midrule
0 & 2 & 2 & 4 & 1{,}5 & 0{,}25\\
1 & 1{,}5 & 0{,}25 & 3 & 1{,}416667 & 0{,}006944\\
2 & 1{,}416667 & 0{,}006944 & 2{,}833333 & 1{,}414216 & $6\times10^{-6}$\\
\bottomrule
\end{tabular}
\end{center}
Em três iterações já temos $\sqrt2$ com 5 casas corretas.
\end{exemplo}

\begin{figure}[h]
\centering
\begin{tikzpicture}
\begin{axis}[width=10cm,height=6cm,axis lines=middle,xmin=1.2,xmax=2.15,ymin=-0.6,ymax=2.3,
  grid=both,grid style={gray!20},xlabel={$x$},ylabel={$y$},xtick={1.416667,1.5,2},
  xticklabels={$x_2$,$x_1$,$x_0$}]
\addplot[domain=1.2:2.1,samples=80,very thick,azul]{x^2-2};
\addplot[domain=1.45:2.1,thick,vermelho]{2+4*(x-2)};
\addplot[domain=1.38:1.75,thick,laranja]{0.25+3*(x-1.5)};
\addplot[dashed] coordinates {(2,0)(2,2)};
\addplot[dashed] coordinates {(1.5,0)(1.5,0.25)};
\addplot[only marks,mark=*] coordinates {(2,2)(1.5,0.25)(1.5,0)(1.416667,0)};
\end{axis}
\end{tikzpicture}
\caption{Newton para $f(x)=x^2-2$: cada nova aproximação é o zero da reta tangente.}
\end{figure}

\paragraph{Cuidados.} Newton exige $f'$ e um bom chute inicial. Se $f'(x_k)\approx0$, o passo
$f(x_k)/f'(x_k)$ fica enorme e a sequência pode ``saltar'' para longe, convergir para outra raiz
ou nem convergir.

\begin{atividade}{Newton-Raphson}{\arqNewton}
A função é $f(x)=x^5-x^3+x^2-3$, e o aplicativo calcula $f'$ e desenha a reta tangente $h$ em
cada passo. O botão INICIAR usa $x_0=1$; o chute também pode ser digitado na caixa $x_0$.
\begin{enumerate}[nosep]
  \item Confira $x_1=1{,}5$, $x_2\approx1{,}339130$, $x_3\approx1{,}286935$, $x_4\approx1{,}282090$,
        $x_5\approx1{,}282051$. Quantas casas corretas cada iteração acrescenta?
  \item Note que $x_1=1{,}5$ passou do outro lado da raiz. Por que, depois disso, todas as
        aproximações se aproximam pela direita? (Observe a concavidade de $f$.)
  \item Experimente $x_0=0$: a tangente é horizontal ($f'(0)=0$) e o método não pode ser aplicado.
        Com $x_0=0{,}3$, $f'(0{,}3)\approx0{,}34$ é pequeno e $x_1\approx8{,}22$; quantas iterações são
        necessárias para voltar perto da raiz?
  \item O botão JAVASCRIPT resolve $f(x)=x^3-x^2-x=0$ com $x_0=0{,}9$ e $\varepsilon=10^{-3}$.
        Como $f'(0{,}9)=-0{,}37$ é pequeno, $x_1\approx-1{,}751$, e o método converge para a raiz
        $\frac{1-\sqrt5}{2}\approx-0{,}618$, e não para a raiz mais próxima,
        $\frac{1+\sqrt5}{2}\approx1{,}618$. Explique o salto com a tangente.
\end{enumerate}
\end{atividade}

% ==================================================================
\section{Método da secante}

Para evitar o cálculo de $f'$, substituímos a tangente pela secante que passa pelos dois últimos
pontos:
\[
x_{k+1}=x_k-\frac{f(x_k)\,(x_k-x_{k-1})}{f(x_k)-f(x_{k-1})}.
\]
São necessários dois valores iniciais $x_0,x_1$. A fórmula lembra a da posição falsa, mas aqui
\emph{não} se mantém um intervalo com troca de sinal: usam-se sempre os dois últimos pontos. A
convergência é superlinear (ordem $\approx1{,}618$), mais rápida que a da posição falsa e um pouco
mais lenta que a de Newton, sem precisar de derivada.

\begin{exemplo}[{: três iterações para $f(x)=x^2-2$ com $x_0=1$, $x_1=2$}]
\begin{center}
\begin{tabular}{cccccc}
\toprule
iteração & $x_{k-1}$ & $x_k$ & $f(x_{k-1})$ & $f(x_k)$ & $x_{k+1}$\\
\midrule
1 & 1 & 2 & $-1$ & 2 & $x_2=4/3\approx1{,}333333$\\
2 & 2 & 1{,}333333 & 2 & $-0{,}222222$ & $x_3=1{,}4$\\
3 & 1{,}333333 & 1{,}4 & $-0{,}222222$ & $-0{,}04$ & $x_4\approx1{,}414634$\\
\bottomrule
\end{tabular}
\end{center}
Na terceira iteração, $f(x_3)$ e $f(x_4)$ têm o mesmo sinal, e mesmo assim o método continua:
não há intervalo a manter.
\end{exemplo}

\begin{atividade}{secante}{\arqSec}
A função é $f(x)=x^2-2$, com $x_0=a=0{,}5$ e $x_1=b=2$ no botão INICIAR. A cada clique em Xk,
o aplicativo traça a secante e substitui $(a,b)$ por $(b,x_{\text{novo}})$.
\begin{enumerate}[nosep]
  \item Confira $x_2=1{,}2$, $x_3=1{,}375$, $x_4\approx1{,}417476$, $x_5\approx1{,}414168$,
        $x_6\approx1{,}414214$.
  \item Em $x_2=1{,}2$ e $x_3=1{,}375$, os dois valores de $f$ são negativos. A posição falsa
        aceitaria esse par? O que isso mostra sobre a diferença entre os dois métodos?
  \item O botão JAVASCRIPT executa o mesmo exemplo com $\varepsilon=10^{-4}$. Compare o número de
        iterações com o da posição falsa no mesmo intervalo (botão JAVASCRIPT do aplicativo de
        posição falsa).
\end{enumerate}
\end{atividade}

% ==================================================================
\section{Comparação dos métodos}

Para $f(x)=x^2-2$, com o critério $|x_k-x_{k-1}|<10^{-6}$ (ou $b_k-a_k<10^{-6}$ na bisseção):

\begin{center}
\begin{tabular}{>{\raggedright\arraybackslash}p{2.6cm}>{\raggedright\arraybackslash}p{3.3cm}c>{\raggedright\arraybackslash}p{2.6cm}>{\raggedright\arraybackslash}p{3.4cm}}
\toprule
Método & Dados iniciais & Iterações & Ordem & Exige\\
\midrule
Bisseção & $[1,2]$ & 20 & linear (fator $\frac12$) & troca de sinal\\
Posição falsa & $[1,2]$ & 9 & linear & troca de sinal\\
Ponto fixo & $x_0=1$, $\varphi=\frac{x+2}{x+1}$ & 9 & linear (fator $|\varphi'(\bar x)|$) & $|\varphi'|<1$ perto da raiz\\
Newton & $x_0=2$ & 5 & quadrática & $f'$ e bom chute\\
Secante & $x_0=1$, $x_1=2$ & 6 & $\approx1{,}618$ & dois chutes\\
\bottomrule
\end{tabular}
\end{center}

\begin{itemize}[nosep]
  \item \textbf{Bisseção}: sempre converge (com troca de sinal) e o número de iterações é
        previsível, mas é lenta. Boa para obter um bom chute inicial.
  \item \textbf{Posição falsa}: em geral mais rápida que a bisseção, mas pode ficar lenta quando
        uma extremidade fica fixa (como no aplicativo com $(x-1)(x-2)^2$).
  \item \textbf{Ponto fixo}: simples, mas a convergência depende da escolha de $\varphi$.
  \item \textbf{Newton}: o mais rápido perto de raízes simples; exige a derivada e pode falhar com
        um chute ruim.
  \item \textbf{Secante}: quase tão rápida quanto Newton, sem derivada; também pode falhar longe da
        raiz.
\end{itemize}
Na prática, é comum combinar: alguns passos de bisseção para chegar perto da raiz e depois
Newton ou secante para refinar.



% ==================================================================
\section{Exercícios}

\begin{enumerate}[label=\textbf{\arabic*.},leftmargin=*]
\item Isole as raízes de $f(x)=x^3-9x+3$ por uma tabela de sinais e justifique que há exatamente
  uma raiz em cada intervalo encontrado.

\item Para $f(x)=e^{-x}-x$ em $[0,1]$:
  \begin{enumerate}[label=(\alph*),nosep]
  \item faça três iterações da bisseção, justificando a escolha de cada intervalo;
  \item quantas iterações garantem $b_k-a_k<10^{-4}$?
  \item faça três iterações da posição falsa e compare.
  \end{enumerate}

\item Aplique o método de Newton a $f(x)=\cos x-x$ com $x_0=1$ até que
  $|x_{k+1}-x_k|<10^{-6}$.

\item Aplique o método da secante a $f(x)=e^{-x}-x$ com $x_0=0$ e $x_1=1$ (quatro iterações).

\item A equação $x^3-x-1=0$ tem uma raiz em $[1,2]$. Considere
  $\varphi_1(x)=\sqrt[3]{x+1}$ e $\varphi_2(x)=x^3-1$.
  \begin{enumerate}[label=(\alph*),nosep]
  \item mostre que as duas são funções de iteração para essa equação;
  \item calcule $|\varphi_1'|$ e $|\varphi_2'|$ perto da raiz e decida qual converge;
  \item faça cinco iterações com a função adequada, a partir de $x_0=1$.
  \end{enumerate}

\item Escreva, em Python, um laço \texttt{while} para o método de Newton aplicado a
  $f(x)=x^2-2$, que pare quando $|x_{k+1}-x_k|<10^{-8}$ ou após 50 iterações. Desenhe o fluxograma
  correspondente.

\item \textbf{(GeoGebra)} No aplicativo de bisseção (versão 1), troque o critério de parada:
  use $tol_1=0{,}1$ e $tol_2=10^{-6}$. Em que iteração o método para? Qual dos dois critérios foi
  decisivo? Explique a diferença entre parar quando \emph{os dois} critérios são satisfeitos e quando
  \emph{um} deles é satisfeito.

\item \textbf{(GeoGebra)} Usando os aplicativos de ponto fixo, proponha uma terceira função de
  iteração para $x^2-x-2=0$, por exemplo $\varphi(x)=\sqrt{x+2}$, e preveja, calculando
  $\varphi'(2)$, se ela converge e se a convergência é monótona ou alternada.
\end{enumerate}

\subsection*{Respostas}
\begin{enumerate}[label=\textbf{\arabic*.},leftmargin=*,nosep]
\item Trocas de sinal em $[-4,-3]$, $[0,1]$ e $[2,3]$; como $f$ tem grau 3, cada intervalo tem
  exatamente uma raiz (também se verifica pelo sinal de $f'(x)=3x^2-9$ em cada um).
\item (a) $x_1=0{,}5$ ($f>0$, novo intervalo $[0{,}5;1]$); $x_2=0{,}75$ ($f<0$, $[0{,}5;0{,}75]$);
  $x_3=0{,}625$ ($f<0$, $[0{,}5;0{,}625]$). (b) $k>\log_2 10^4\approx13{,}3$: 14 iterações.
  (c) $x_1\approx0{,}612700$, $x_2\approx0{,}572181$, $x_3\approx0{,}567703$ (com $a=0$ fixo);
  a raiz é $\bar x\approx0{,}567143$.
\item $x_1\approx0{,}750364$, $x_2\approx0{,}739113$, $x_3\approx0{,}739085133$,
  $x_4\approx0{,}739085133$.
\item $x_2\approx0{,}612700$, $x_3\approx0{,}563838$, $x_4\approx0{,}567170$, $x_5\approx0{,}567143$.
\item (b) Na raiz $\bar x\approx1{,}3247$: $|\varphi_1'(\bar x)|\approx0{,}19<1$ (converge) e
  $|\varphi_2'(\bar x)|=3\bar x^2\approx5{,}3>1$ (diverge). (c) $1{,}259921$; $1{,}312294$;
  $1{,}322354$; $1{,}324269$; $1{,}324633$.
\item Uma solução:
\begin{lstlisting}[language=Python]
x, k = 2.0, 0
while k < 50:
    x_novo = x - (x**2 - 2) / (2*x)
    k = k + 1
    if abs(x_novo - x) < 1e-8:
        break
    x = x_novo
print(k, x_novo)
\end{lstlisting}
\item Para em $k=15$ ($x_k\approx1{,}7632233$, $|f(x_k)|\approx7\times10^{-7}$). O critério
  do intervalo ($|b-a|\le0{,}1$) é atingido já em $k=1$; o método só para quando também
  $|f(x_k)|\le10^{-6}$, portanto é o segundo critério que decide. Parar com \emph{um} critério
  economiza iterações, mas pode parar cedo demais.
\item $\varphi'(x)=\frac{1}{2\sqrt{x+2}}$, $\varphi'(2)=\frac14$: converge, de forma monótona
  (pois $\varphi'>0$).
\end{enumerate}
